% Responsable: Danela
\subsection{Sistema de un GAD con contrato público}

\subsubsection{Caracterización}
El docente presentó el escenario sin detalle, por lo que se describen los supuestos usados para el análisis.

\begin{table}[H]
\centering\footnotesize
\caption{Caracterización del escenario del GAD}\label{tab:gad}
\begin{tabularx}{\textwidth}{@{}p{0.26\textwidth}Y@{}}
\toprule
\textbf{Aspecto} & \textbf{Descripción del escenario}\\
\midrule
Sistema & Se asume que es una plataforma web para que los ciudadanos realicen trámites municipales en línea, ya sean solicitudes, pagos de tasas y seguimiento de trámites, para que los funcionarios los atiendan.\\
Cliente & El GAD municipal, representado por un administrador del contrato designado por la institución.\\
Actores del proceso & Administrador del contrato y fiscalizador por parte del GAD, funcionarios de los departamentos que atienden trámites, ciudadanos usuarios, autoridades del GAD que aprueban el alcance, el equipo proveedor y los organismos de control del uso de fondos públicos.\\
Estabilidad de los requisitos (C1) & Alta, ya que se asume que los requisitos se fijan en los términos de referencia del proceso de contratación pública antes de firmar el contrato, y cualquier cambio requiere una modificación formal.\\
Criticidad (C2) & Media, porque una falla no pone en riesgo vidas, pero afecta la atención a los ciudadanos y el manejo de datos personales y pagos.\\
Disponibilidad del cliente (C3) & Baja, ya que los funcionarios tienen otras responsabilidades y las decisiones pasan por aprobaciones internas, por lo que las reuniones son programadas y poco frecuentes.\\
Tamaño y complejidad (C4) & Grande, porque involucra varios departamentos, muchos tipos de trámite y la integración con sistemas existentes del GAD.\\
Regulación y contrato (C5) & Alta, el proyecto se rige por un contrato público con alcance, plazo y presupuesto definidos, y debe cumplir la normativa de contratación pública y de protección de datos personales.\\
Documentación exigida (C6) & Alta, porque se necesitan especificaciones, actas de entrega y recepción e informes de fiscalización para pagar cada entregable y responder ante los organismos de control.\\
\bottomrule
\end{tabularx}
\end{table}

\subsubsection{Modelo elegido y justificación}
\textbf{Modelo elegido:} Cascada

Se elige el modelo en cascada por las siguientes razones del escenario:

\begin{enumerate}
\item Los requisitos quedan fijados en los términos de referencia antes de firmar el contrato, y cualquier cambio exige una modificación formal. La cascada parte justamente de requisitos estables definidos al inicio.
\item El contrato público define alcance, plazo y presupuesto cerrados, y los pagos se hacen contra entregables aprobados por el fiscalizador. Las fases secuenciales de la cascada, con un entregable revisado al final de cada una, se ajustan a esa forma de aceptar el producto.
\item La disponibilidad del cliente es baja y la cascada concentra la participación de los funcionarios en la obtención y validación inicial de los requisitos, y no depende de reuniones frecuentes en cada iteración.
\item Se exige mucha documentación formal para pagar cada entregable y responder ante los organismos de control, la cascada produce una especificación completa y documentos de cada fase que sirven como evidencia.
\end{enumerate}

\textbf{Modelos descartados:}

\textbf{Ágil:} Supone requisitos cambiantes y un cliente disponible en cada iteración para priorizar y validar, condiciones que no se cumplen en el GAD. Además, su documentación mínima no alcanza para las actas y los informes que exige el contrato.

\textbf{Iterativo:} Se adapta mejor al tamaño del proyecto, porque permite entregar por partes la integración de varios departamentos. Sin embargo, revisar los requisitos en cada iteración choca con un alcance fijado por contrato y con un cliente poco disponible, por lo que obtiene menor total que la cascada en la matriz.

\subsubsection{Riesgos del modelo elegido y mitigación}
% Ronda 3 (sábado, hasta 22:30): riesgos concretos y cómo mitigarlos.
